\section{Estructura de la memoria}
La presente memoria se estructura en los siguientes capítulos, diseñados para guiar al lector de forma lógica a través del desarrollo del proyecto:

\begin{itemize}
    
    \item En el \textbf{\autoref{ch:estado_arte}}, se presenta el estado del arte y el marco teórico que fundamenta la investigación. Se comienza con una contextualización clínica de las enfermedades neurodegenerativas con mayor prevalencia (Parkinson, Esclerosis Lateral Amiotrófica y Alzheimer). A continuación, se introduce el uso de la voz como biomarcador y se detalla el flujo de trabajo técnico necesario para su análisis. Posteriormente, se realiza una revisión exhaustiva y una comparativa de los modelos de Inteligencia Artificial empleados en la literatura, analizando sus ventajas y limitaciones, y finalizando con una exploración de las tendencias futuras.
    
    \item En el \textbf{\autoref{ch:metodologia}}, se describen los recursos materiales y técnicos empleados. Se detallan los conjuntos de datos (datasets) utilizados y el entorno de desarrollo. Asimismo, se expone exhaustivamente el pipeline de preprocesamiento aplicado a las señales de audio, tanto para la fase de entrenamiento y validación como para el entorno de inferencia. Por último, se definen las métricas de evaluación que determinarán el rendimiento y la validez de los modelos.
    
    \item En el \textbf{\autoref{ch:diseno}}, se profundiza en la arquitectura e implementación de los modelos evaluados. Se detallan las configuraciones experimentales y las estrategias de entrenamiento, tales como la optimización de hiperparámetros y el uso de validación cruzada. El capítulo abarca desde el modelado con Machine Learning tradicional hasta arquitecturas de Deep Learning y el ajuste fino (Fine-Tuning) de Foundation Models.
    
    \item En el \textbf{\autoref{ch:prueba_concepto}}, se expone el diseño y desarrollo de la aplicación cliente-servidor construida como prueba de concepto. Se detalla la arquitectura de software elegida, describiendo los componentes del sistema (Frontend y API REST Backend), el flujo de datos durante la inferencia y el despliegue del modelo final junto con la generación de explicabilidad (SHAP), concluyendo con el diseño de la interfaz de usuario.
    
    \item En el \textbf{\autoref{ch:resultados}}, se presentan, evalúan y discuten los resultados obtenidos. Se realiza una comparativa del rendimiento de los modelos bajo distintos criterios, incluyendo la modalidad de los datos, la complejidad arquitectónica y la interpretabilidad clínica, culminando con un análisis global del compromiso entre precisión y coste computacional.
    
    \item En el \textbf{\autoref{ch:conclusion}}, se sintetizan las conclusiones finales, evaluando el grado de consecución de los objetivos planteados al inicio del proyecto. Además, se identifican las limitaciones encontradas durante el desarrollo y se proponen posibles líneas de trabajo futuro para dar continuidad a la investigación.
    
\end{itemize} 

Adicionalmente, se incluye material complementario al final del documento estructurado en los siguientes anexos:

\begin{itemize}
    \item En el \textbf{\hyperref[app:anexo_eda]{Apéndice~\ref*{app:anexo_eda}}}, se proporciona todo el análisis exploratorio de los datasets de forma extensa, se muestran todo tipo de informacion y distribución de los atributos demográficos de los pacientes registrados y de las variables acústicas extraídas de los audios.
    \item En el \textbf{\hyperref[app:anexo_entorno]{Apéndice~\ref*{app:anexo_entorno}}}, se proporciona el listado completo de las dependencias del entorno virtual empleado, asegurando la reproducibilidad del entorno de desarrollo.
    \item En el \textbf{\hyperref[app:anexo_repositorio]{Apéndice~\ref*{app:anexo_repositorio}}}, se detalla la ubicación del repositorio público de código fuente, así como los enlaces y credenciales (mediante códigos QR) para acceder a los binarios compilados de la aplicación móvil (instalables para plataformas Android e iOS).
\end{itemize}